% Build with: latexmk -pdf main.tex
\documentclass[11pt]{article}

\usepackage[T1]{fontenc}
\usepackage{lmodern}
\usepackage[margin=1in]{geometry}
\usepackage{microtype}
\usepackage{amsmath,amssymb,amsthm,mathtools}
\usepackage{booktabs}
\usepackage{enumitem}
\usepackage{xcolor}
\usepackage[colorlinks=true,linkcolor=blue!50!black,citecolor=blue!50!black,urlcolor=blue!50!black]{hyperref}

\theoremstyle{definition}
\newtheorem{definition}{Definition}[section]
\theoremstyle{plain}
\theoremstyle{remark}
\newtheorem{remark}[definition]{Remark}

\newcommand{\bits}[1]{\{0,1\}^{#1}}
\newcommand{\getsr}{\stackrel{\$}{\gets}}
\newcommand{\Th}{\mathsf{Th}}
\newcommand{\Enc}{\mathsf{Enc}}
\newcommand{\Digest}{\mathsf{Digest}}
\newcommand{\Gen}{\mathsf{Gen}}
\newcommand{\Sig}{\mathsf{Sig}}
\newcommand{\Ver}{\mathsf{Ver}}
\newcommand{\SIG}{\mathsf{SIG}}
\newcommand{\Chain}{\mathsf{Chain}}
\newcommand{\hash}{\mathsf{H}}
\newcommand{\LE}{\mathsf{LE}}
\newcommand{\Truncate}{\mathsf{Truncate}}
\newcommand{\concat}{\mathbin\Vert}
\newcommand{\sk}{\mathit{sk}}
\newcommand{\pk}{\mathit{pk}}
\newcommand{\rootnode}{\mathit{root}}
\newcommand{\tw}{\mathit{tw}}
\newcommand{\lmsg}{\ell_{\mathrm{msg}}}
\newcommand{\lpar}{\ell_{\mathrm{p}}}
\newcommand{\ltwk}{\ell_{\mathrm{t}}}
\newcommand{\lrnd}{\ell_{\mathrm{rnd}}}
\newcommand{\lctr}{\ell_{\mathrm{c}}}
\newcommand{\qs}{q_{\mathrm{s}}}
\newcommand{\amax}{A_{\max}}
\newcommand{\cmax}{C_{\max}}
\newcommand{\idx}{\mathit{idx}}
\newcommand{\lay}{\mathit{lay}}
\newcommand{\OtsSign}{\mathsf{Ots.sign}}
\newcommand{\OtsLeaf}{\mathsf{Ots.leaf}}
\newcommand{\TreeRoot}{\mathsf{Tree.root}}
\newcommand{\TreePath}{\mathsf{Tree.path}}
\newcommand{\TreeFold}{\mathsf{Tree.fold}}
\newcommand{\FtsKey}{\mathsf{Fts.key}}
\newcommand{\FtsOpen}{\mathsf{Fts.open}}
\newcommand{\FtsRec}{\mathsf{Fts.recover}}

\emergencystretch=1.5em

\title{Example of a SPHINCS$^+$ variant}
\author{}
\date{}

\begin{document}
\maketitle

\begin{abstract}

We present, as an example, a SPHINCS$^+$-based signature with the following properties:

\begin{itemize}
  \item \textbf{stateless}: supporting up to $2^{24}$ signatures.
  \item \textbf{NIST security level~1}~\cite{NISTPQC} (TODO prove it)
  \item \textbf{public key: 32 bytes}.
  \item \textbf{signature: 4924 bytes}.
  \item \textbf{497 hashes per verification}.
  \item signing costs 190K hashes with 1024 bytes of cached signer state, or 1.55M without.
  \item \textbf{key generation costs 1.38M hashes}.
\end{itemize}
\end{abstract}

The construction is SPHINCS$^+$~\cite{SPHINCSPLUS,FIPS205} with two of the optimizations surveyed in~\cite{KN25}, WOTS$^+$C and FORS$^+$C, both from~\cite{HK22C}; its third, PORS$^+$FP, is not used.

\section{Definitions and notation}

\begin{definition}[Signature scheme]
A signature scheme is a tuple $\SIG=(\Gen,\Sig,\Ver)$, where $\Gen$ and $\Sig$ are randomized and $\Ver$ is deterministic:
\[
  \Gen\longrightarrow(\pk,\sk),\qquad
  \Sig(\sk,m)\longrightarrow\sigma\in\Sigma\cup\{\bot\},\qquad
  \Ver(\pk,m,\sigma)\longrightarrow\{0,1\},
\]
where $\Sigma$ is the signature space and $m\in\bits{\lmsg}$. Whenever $(\pk,\sk)$ is output by $\Gen$ and $\Sig(\sk,m)$ returns $\sigma\neq\bot$, correctness requires $\Ver(\pk,m,\sigma)=1$. $\Sig$ keeps no state and may be called on any message any number of times, but security degrades with that number: this specification is stated for at most $\qs$ signatures per key pair.
\end{definition}

Byte strings are concatenated with $\concat$. Bits and integer encodings are little endian. $\LE_r(a)$ is the unsigned $r$-bit encoding of $a$. All indices are zero based. Layers are numbered from the top: layer $0$ carries the public key, layer $d-1$ signs few-time keys.

\begin{center}
\begin{tabular}{@{}lll@{}}
\toprule
Symbol & Value & Meaning\\
\midrule
$n$ & $128$ bits & hash value and Merkle node length\\
$\lpar$ & $128$ bits & public parameter length\\
$\ltwk$ & $128$ bits & tweak length\\
$\lmsg$ & $256$ bits & message length\\
$\lrnd$ & $128$ bits & randomizer length\\
$\lctr$ & $32$ bits & encoding counter length\\
$w$ & $3$ & chunk size in bits\\
$v$ & $42$ & code length\\
$T$ & $191$ & target sum\\
$d$ & $3$ & hypertree layers\\
$(h_0,h_1,h_2)$ & $(12,7,7)$ & Merkle tree height of each layer\\
$h$ & $26$ & total height, $h=\sum_\lay h_\lay$\\
$a$ & $10$ & $\log_2$ of the leaves in one few-time tree\\
$k$ & $15$ & digest index groups; the forest holds $k-1$ trees\\
$\qs$ & $2^{24}$ & signatures per key pair\\
$\amax$ & $2^{32}$ & maximum digest attempts per signature\\
$\cmax$ & $2^{32}$ & maximum encoding attempts per layer\\
\bottomrule
\end{tabular}
\end{center}

Let $\hash:\bits{*}\to\bits{256}$ be a cryptographic hash function, and let $\Truncate_\nu$ keep the first $\nu$ bits of its output. The tweakable hash $\Th:\mathcal P\times\mathcal T\times\mathcal M\to\mathcal H$, with $\mathcal P=\bits{\lpar}$, $\mathcal T=\bits{\ltwk}$, $\mathcal M=\bits{*}$ and $\mathcal H=\bits{n}$, is
\[
  \Th(P,\tw,M)=\Truncate_n\!\left(\hash(\tw\concat P\concat M)\right).
\]
$\hash$ and the code $\mathcal C$ of Section~\ref{sec:ots} are those of~\cite{leanVMb}, with a different target sum.

\begin{definition}[Tweak encoding]
For one-byte $t$ and $\lay$, and unsigned 32-bit integers $\tau$, $p$ and $j$, define the 16-byte tweak
\[
  \mathsf{enc}(t,\lay,\tau,p,j)=\LE_8(t)\concat\LE_8(\lay)\concat\LE_{32}(\tau)\concat\LE_{32}(p)\concat\LE_{32}(j)\concat\LE_{16}(0),
\]
fourteen bytes of fields and two of padding. The byte-wide fields cap $d\leq256$ and $k\leq257$; the 32-bit fields are never near their range here.
\end{definition}

A tweak names one hash call in the whole structure, which is what lets a security argument treat each call separately. Inside the hypertree, $\lay$ is the layer and $\tau$ the tree within it; inside a few-time key, $\lay$ is the tree in the forest and $\tau$ the index $\idx$ that selects the instance. Define
\[
\begin{aligned}
  \mathsf{tw}_{\mathrm{prf}}(\lay,\tau,i,e) &= \mathsf{enc}(0,\lay,\tau,i,e),\\
  \mathsf{tw}_{\mathrm{chain}}(\lay,\tau,e,i,\mu) &= \mathsf{enc}(1,\lay,\tau,2^wi+\mu-1,e), &&0\leq i<v,\ 0<\mu<2^w,\\
  \mathsf{tw}_{\mathrm{leaf}}(\lay,\tau,e) &= \mathsf{enc}(2,\lay,\tau,0,e),\\
  \mathsf{tw}_{\mathrm{node}}(\lay,\tau,\lambda,j) &= \mathsf{enc}(3,\lay,\tau,\lambda,j), &&1\leq\lambda\leq h_\lay,\ j<2^{h_\lay-\lambda},\\
  \mathsf{tw}_{\mathrm{enc}}(\lay,\tau,e) &= \mathsf{enc}(4,\lay,\tau,0,e),\\
  \mathsf{tw}_{\mathrm{ftsprf}}(\idx,\kappa,j) &= \mathsf{enc}(5,\kappa,\idx,0,j),\\
  \mathsf{tw}_{\mathrm{ftsleaf}}(\idx,\kappa,j) &= \mathsf{enc}(6,\kappa,\idx,0,j),\\
  \mathsf{tw}_{\mathrm{ftsnode}}(\idx,\kappa,\lambda,j) &= \mathsf{enc}(7,\kappa,\idx,\lambda,j), &&1\leq\lambda\leq a,\ j<2^{a-\lambda},\\
  \mathsf{tw}_{\mathrm{ftsroots}}(\idx) &= \mathsf{enc}(8,0,\idx,0,0),\\
  \mathsf{tw}_{\mathrm{msg}} &= \mathsf{enc}(9,0,0,0,0).
\end{aligned}
\]
Every structural position has its own tweak. Two are shared on purpose: $\mathsf{tw}_{\mathrm{msg}}$ across every message and $\mathsf{tw}_{\mathrm{enc}}(\lay,\tau,e)$ across every counter tried at that leaf. Types $0$ and $5$ are used only by the seed derivation of Remark~\ref{rem:seed}.

\section{The index}

An index $\idx\in\{0,\ldots,2^h-1\}$ says which few-time key signs, and with it which tree and which leaf are used on every layer:
\[
  \tau_\lay(\idx)=\left\lfloor\idx\,/\,2^{\sum_{j\geq\lay}h_j}\right\rfloor,
  \qquad
  e_\lay(\idx)=\left\lfloor\idx\,/\,2^{\sum_{j>\lay}h_j}\right\rfloor\bmod 2^{h_\lay}.
\]
With $(h_0,h_1,h_2)=(12,7,7)$ the divisors are $2^{26},2^{14},2^{7}$ for $\tau$ and $2^{14},2^{7},2^{0}$ for $e$. Layer $0$ has $\tau_0=0$, its single tree being the public key, and layer $d-1$ has $e_{d-1}=\idx\bmod2^{h_{d-1}}$. The layers link through the same two functions,
\[
  \tau_\lay=\tau_{\lay-1}\cdot2^{h_{\lay-1}}+e_{\lay-1},
\]
so the tree used on layer $\lay$ is the one whose root sits at leaf $e_{\lay-1}$ of the tree used on layer $\lay-1$. Layer $\lay$ holds $2^{\sum_{j<\lay}h_j}$ trees of $2^{h_\lay}$ leaves, so $(\tau_\lay,e_\lay)$ takes $2^{\sum_{j\leq\lay}h_j}$ values, that is $2^{12}$, $2^{19}$ and $2^{26}$ here, the last putting the $2^h$ indices in bijection with the leaves of the bottom layer.

\section{The one-time signature}
\label{sec:ots}

One position $(\lay,\tau,e)$ holds one one-time key: $v$ secret values, one per hash chain of $2^w-1$ steps. A signature on $M$ encodes $M$ into a codeword $x$ and reveals, on each chain, the value at position $x_i$; a verifier walks the remaining $2^w-1-x_i$ steps and recovers the public value. Forging on $M'\neq M$ needs a codeword $x'$ with $x'_i\geq x_i$ everywhere, so that every revealed value lies below what the forger needs, and otherwise a chain must be inverted. The codewords are therefore chosen to make that impossible:
\[
  \mathcal C=\left\{(x_0,\ldots,x_{v-1})\in\{0,\ldots,2^w - 1\}^{v}:\sum_{i=0}^{v-1}x_i=T\right\}.
\]
Two distinct words of equal sum cannot be ordered componentwise, so $x'\geq x$ forces $x'=x$: the code is incomparable. This is what removes the Winternitz checksum of the classical scheme, since the verifier checks the sum itself, and it is why a counter is needed, most messages not encoding into $\mathcal C$ at all. That counter is carried in the signature: WOTS$^+$C~\cite{HK22C}.

\begin{definition}[Encoding]
$\Enc(P,\lay,\tau,e,M,c)$ computes
\[
  D=\Th\!\left(P,\mathsf{tw}_{\mathrm{enc}}(\lay,\tau,e),M\concat\LE_{\lctr}(c)\right),
\]
writes $D=D_0\concat D_1$ with $D_0,D_1$ of $n/2$ bits, and lets $d_q$ be the integer encoded little endian by $D_q$. For $q\in\{0,1\}$ and $0\leq r<v/2$ it sets
\[
  x_{qv/2+r}=\left\lfloor\frac{d_q}{2^{wr}}\right\rfloor\bmod 2^w .
\]
It returns $x=(x_0,\ldots,x_{v-1})$ if bit $wv/2=63$ of both $d_0$ and $d_1$ is zero and $x\in\mathcal C$; otherwise $\bot$.
\end{definition}

Each half of $D$ carries $v/2=21$ chunks, that is $wv/2=63$ bits, and one pinned bit, so all $2(63+1)=128$ digest bits are accounted for and $x$ determines $D$. A uniform $D$ is admissible with probability $p=2^{-13.60}$, so a signer tries $1/p=12{,}436$ counters per layer on average.

\begin{definition}[Hash chain]
For a chain index $i$ at position $(\lay,\tau,e)$, a start $\mu$, a number of steps $s$ and a value $y\in\mathcal H$, define $\Chain_{\lay,\tau,e,i}(P,\mu,0,y)=y$ and, for $s\geq1$,
\[
  \Chain_{\lay,\tau,e,i}(P,\mu,s,y)=\Th\!\left(P,\mathsf{tw}_{\mathrm{chain}}(\lay,\tau,e,i,\mu+s),\Chain_{\lay,\tau,e,i}(P,\mu,s-1,y)\right),
\]
valid for $0\leq\mu<2^w$ and $0\leq s\leq2^w-1-\mu$.
\end{definition}

The secrets $\sk_{\lay,\tau,e,i}$ for $0\leq i<v$ are sampled by $\Gen$; Remark~\ref{rem:seed} derives them instead. The public values and the Merkle leaf they form are
\[
\begin{aligned}
  \pk_{\lay,\tau,e,i}&=\Chain_{\lay,\tau,e,i}\!\left(P,0,2^w-1,\sk_{\lay,\tau,e,i}\right),\\
  X^{\lay,\tau}_{0,e}&=\Th\!\left(P,\mathsf{tw}_{\mathrm{leaf}}(\lay,\tau,e),\pk_{\lay,\tau,e,0}\concat\cdots\concat\pk_{\lay,\tau,e,v-1}\right).
\end{aligned}
\]

\begin{definition}[$\OtsSign$]
$\OtsSign(P,\lay,\tau,e,M)$ takes the least $c<\cmax$ for which
\[
  x=\Enc(P,\lay,\tau,e,M,c)\neq\bot,
\]
returning $\bot$ if there is none, and outputs $(c,\sigma)$ with
\[
  \sigma_i=\Chain_{\lay,\tau,e,i}\!\left(P,0,x_i,\sk_{\lay,\tau,e,i}\right),\qquad 0\leq i<v.
\]
\end{definition}

\begin{definition}[$\OtsLeaf$]
$\OtsLeaf(P,\lay,\tau,e,M,c,\sigma)$ computes $x=\Enc(P,\lay,\tau,e,M,c)$ and returns $\bot$ if it is $\bot$; otherwise it returns
\[
  \Th\!\left(P,\mathsf{tw}_{\mathrm{leaf}}(\lay,\tau,e),\pk'_0\concat\cdots\concat\pk'_{v-1}\right),
  \qquad
  \pk'_i=\Chain_{\lay,\tau,e,i}\!\left(P,x_i,2^w-1-x_i,\sigma_i\right).
\]
\end{definition}

$\OtsLeaf$ returns $X^{\lay,\tau}_{0,e}$ on what $\OtsSign$ produced for that $M$, and on nothing else short of colliding $\Th$: a second admissible counter for the same codeword would do, and by the previous paragraph that codeword fixes $D$, so it is a $128$-bit hit. Otherwise it returns a value the forger does not control. It is deterministic and does not touch the secrets, which is why it is the verifier's half.

\begin{remark}[One key, one message]
\label{rem:once}
Each one-time key signs exactly one value. The key at leaf $e_\lay$ of tree $(\lay,\tau_\lay)$ is reached only by indices with those two components, and every such index gives it the same $M$, namely the root of the tree below at $\lay<d-1$ and the few-time public key at $\lay=d-1$. This is why $\OtsSign$ fixes the counter as a function of its inputs, here the \emph{least} admissible one, rather than taking any admissible one. A signer that resumed or randomized that search would make one key sign two distinct codewords $x\neq x'$ of equal sum; being incomparable neither covers the other, so the pair leaves an adversary holding $\min(x_i,x'_i)$ on every chain, and forging then costs only grinding $(M',c')$ until $\Enc$ yields a codeword dominating that minimum, of order $2^{53}$ for these parameters against the $2^{128}$ intended.
\end{remark}

\section{A layer}
\label{sec:layer}

Layer $\lay$'s tree $\tau$ is the Merkle tree over the $2^{h_\lay}$ one-time leaves $X^{\lay,\tau}_{0,e}$:
\[
  X^{\lay,\tau}_{\lambda,j}=\Th\!\left(P,\mathsf{tw}_{\mathrm{node}}(\lay,\tau,\lambda,j),X^{\lay,\tau}_{\lambda-1,2j}\concat X^{\lay,\tau}_{\lambda-1,2j+1}\right),
  \qquad 1\leq\lambda\leq h_\lay,\ 0\leq j<2^{h_\lay-\lambda}.
\]

\begin{definition}[$\TreeRoot$, $\TreePath$]
$\TreeRoot(P,\lay,\tau)=X^{\lay,\tau}_{h_\lay,0}$, and $\TreePath(P,\lay,\tau,e)=(A_0,\ldots,A_{h_\lay-1})$ with
\[
  A_\lambda=X^{\lay,\tau}_{\lambda,\lfloor e/2^\lambda\rfloor\oplus1}.
\]
\end{definition}

\begin{definition}[$\TreeFold$]
$\TreeFold(P,\lay,\tau,e,X,A)$ sets $V_0=X$ and, for $0\leq\lambda<h_\lay$ with $j=\lfloor e/2^{\lambda+1}\rfloor$,
\[
  V_{\lambda+1}=
  \begin{cases}
    \Th\!\left(P,\mathsf{tw}_{\mathrm{node}}(\lay,\tau,\lambda+1,j),V_\lambda\concat A_\lambda\right) & \text{if bit $\lambda$ of $e$ is }0,\\[2pt]
    \Th\!\left(P,\mathsf{tw}_{\mathrm{node}}(\lay,\tau,\lambda+1,j),A_\lambda\concat V_\lambda\right) & \text{otherwise,}
  \end{cases}
\]
and returns $V_{h_\lay}$.
\end{definition}

$\TreeFold(P,\lay,\tau,e,X^{\lay,\tau}_{0,e},\TreePath(P,\lay,\tau,e))=\TreeRoot(P,\lay,\tau)$: the path and the fold are the two halves of one Merkle opening.

\section{The few-time signature}
\label{sec:fts}

A few-time key is a forest of $k-1$ Merkle trees of $2^a$ secret leaves. A signature opens one leaf per tree, at indices a digest chooses. Reuse leaks rather than breaks: after $r$ signatures under one instance an adversary holds $r$ leaves per tree, and it can sign any message that lands on that instance, has last index zero, and has its other indices all on leaves it holds, the last condition holding with probability $\left(1-(1-2^{-a})^{r}\right)^{k-1}$. The forest holds $k-1$ trees rather than $k$ because the digest is resampled until its last index is zero, so that tree carries no information and is dropped: FORS$^+$C~\cite{HK22C}.

The secrets $s_{\idx,\kappa,j}$ for $0\leq\kappa<k-1$ and $0\leq j<2^a$ are sampled by $\Gen$; Remark~\ref{rem:seed} derives them instead. Their trees are
\[
\begin{aligned}
  Y^{\idx,\kappa}_{0,j} &= \Th\!\left(P,\mathsf{tw}_{\mathrm{ftsleaf}}(\idx,\kappa,j),s_{\idx,\kappa,j}\right),\\
  Y^{\idx,\kappa}_{\lambda,j} &= \Th\!\left(P,\mathsf{tw}_{\mathrm{ftsnode}}(\idx,\kappa,\lambda,j),Y^{\idx,\kappa}_{\lambda-1,2j}\concat Y^{\idx,\kappa}_{\lambda-1,2j+1}\right), &&1\leq\lambda\leq a.
\end{aligned}
\]

\begin{definition}[$\FtsKey$]
$\FtsKey(P,\idx)=\Th\!\left(P,\mathsf{tw}_{\mathrm{ftsroots}}(\idx),Y^{\idx,0}_{a,0}\concat\cdots\concat Y^{\idx,k-2}_{a,0}\right)$.
\end{definition}

\begin{definition}[$\FtsOpen$]
$\FtsOpen(P,\idx,u)$, for indices $u=(u_0,\ldots,u_{k-1})$, outputs $(s,B)$ with, for $0\leq\kappa<k-1$,
\[
  s_\kappa=s_{\idx,\kappa,u_\kappa},
  \qquad
  B_{\kappa,\lambda}=Y^{\idx,\kappa}_{\lambda,\lfloor u_\kappa/2^\lambda\rfloor\oplus1},\qquad 0\leq\lambda<a.
\]
It ignores $u_{k-1}$, whose tree is the dropped one.
\end{definition}

\begin{definition}[$\FtsRec$]
$\FtsRec(P,\idx,u,s,B)$ sets, for $0\leq\kappa<k-1$, $W_{\kappa,0}=\Th(P,\mathsf{tw}_{\mathrm{ftsleaf}}(\idx,\kappa,u_\kappa),s_\kappa)$ and, for $0\leq\lambda<a$ with $j=\lfloor u_\kappa/2^{\lambda+1}\rfloor$,
\[
  W_{\kappa,\lambda+1}=
  \begin{cases}
    \Th\!\left(P,\mathsf{tw}_{\mathrm{ftsnode}}(\idx,\kappa,\lambda+1,j),W_{\kappa,\lambda}\concat B_{\kappa,\lambda}\right) & \text{if bit $\lambda$ of $u_\kappa$ is }0,\\[2pt]
    \Th\!\left(P,\mathsf{tw}_{\mathrm{ftsnode}}(\idx,\kappa,\lambda+1,j),B_{\kappa,\lambda}\concat W_{\kappa,\lambda}\right) & \text{otherwise,}
  \end{cases}
\]
and returns $\Th\!\left(P,\mathsf{tw}_{\mathrm{ftsroots}}(\idx),W_{0,a}\concat\cdots\concat W_{k-2,a}\right)$.
\end{definition}

$\FtsRec$ returns $\FtsKey(P,\idx)$ on an opening of the leaves $u$ of that instance, and on nothing else short of colliding $\Th$.

\begin{definition}[Message digest]
$\Digest(P,\rootnode,m,\rho)$ returns
\[
  \Truncate_{h+ka}\!\left(\hash\!\left(\mathsf{tw}_{\mathrm{msg}}\concat P\concat \rho\concat\rootnode\concat m\right)\right),
\]
$22$ bytes, that is $h+ka=176$ bits. Let $N$ be the integer those bytes encode little endian. The index and the leaf indices are
\[
  \idx=N\bmod 2^{h},\qquad
  u_\kappa=\left\lfloor N\,/\,2^{h+\kappa a}\right\rfloor\bmod 2^{a},\qquad 0\leq\kappa<k,
\]
and the digest is admissible if $u_{k-1}=0$.
\end{definition}

\begin{remark}[The randomizer]
\label{rem:rho}
$\rho\in\bits{\lrnd}$ is sampled fresh per attempt and carried in the signature, the verifier recomputing the digest from it. Without it the argument would need collision resistance of $\Digest$: an adversary could look for two messages sharing a digest and carry a signature from either to the other. Sampling denies it that, since it cannot know, before committing to a message, which digest that message will be signed under. FIPS~205 has the same element as $R$, derives it from a secret key and an optional input rather than sampling it, and defaults to drawing that input at random, keeping a deterministic choice for platforms without a random source~\cite{FIPS205}.

Resampling until $u_{k-1}=0$ is what gives FORS$^+$C its dropped tree. It conditions the digest without biasing $\idx$, the $h$ index bits and the $a$ bits of $u_{k-1}$ being disjoint parts of a random oracle output.
\end{remark}

\section{Key generation}

$\Gen$ independently samples $P\getsr\bits{\lpar}$, the one-time secrets $\sk_{\lay,\tau,e,i}\getsr\mathcal H$ for every layer, tree, leaf and chain, and the few-time secrets $s_{\idx,\kappa,j}\getsr\mathcal H$ for every instance, tree and leaf. It returns
\[
  \pk=(\rootnode,P),\qquad \rootnode=\TreeRoot(P,0,0),\qquad
  \sk=\left(P,\rootnode,(\sk_{\lay,\tau,e,i}),(s_{\idx,\kappa,j})\right),
\]
with $\pk$ serialized as $\rootnode\concat P$, 32 bytes. $\rootnode$ is in $\sk$ because $\Sig$ hashes it into every digest; it is not a secret and is recomputable. Only layer $0$ is built here; the trees below it are built when a signature needs them. Restricting a key to a range of indices is not available, $\idx$ being derived from the message, so a key answers for all $2^h$ of them.

\begin{remark}[Seed derivation]
\label{rem:seed}
A key pair as specified holds one secret per chain and per few-time leaf, which no signer can store. They may instead be generated from a single master secret $S\getsr\mathcal H$ by
\[
  \sk_{\lay,\tau,e,i}=\Th\!\left(P,\mathsf{tw}_{\mathrm{prf}}(\lay,\tau,i,e),S\right),
  \qquad
  s_{\idx,\kappa,j}=\Th\!\left(P,\mathsf{tw}_{\mathrm{ftsprf}}(\idx,\kappa,j),S\right),
\]
leaving $\sk=(P,\rootnode,S)$, 48 bytes. The specification samples rather than derives because the security statement is about independent uniform secrets and a derivation is an implementation of it, sound in the ROM; every cost in Section~\ref{sec:costs} assumes the derivation, since the alternative is unimplementable.
\end{remark}

\begin{remark}[Signer state]
\label{rem:cache}
Layer $0$'s tree is the same for every signature, so a signer may keep part of it. Keeping the $2^6$ nodes at depth $\lceil h_0/2\rceil=6$, 1024 bytes, leaves a $2^6$-leaf subtree to rebuild per signature and at most $2^6-1$ nodes to refold above it: that layer's share of signing falls from 1.38M hashes to 21.7K, and the total from 1.55M to 190K. It is a cache, not state: a deterministic function of the secrets, so losing it costs recomputation. Traversal~\cite{BDS08} would amortize further, but consecutive signatures land on unrelated leaves.
\end{remark}

\section{Signing}
\label{sec:sign}

$\Sig(\sk,m)$:
\begin{enumerate}[leftmargin=2em]
\item Make at most $\amax$ attempts; in each, sample $\rho\getsr\bits{\lrnd}$ independently and compute $\Digest(P,\rootnode,m,\rho)$, stopping at the first admissible digest, which yields $\idx$ and $u$. Return $\bot$ if every attempt fails. This takes $2^a=1024$ attempts on average.
\item $(s,B)\gets\FtsOpen(P,\idx,u)$ and $M_{d-1}\gets\FtsKey(P,\idx)$.
\item For $\lay=d-1$ down to $0$, with $\tau=\tau_\lay(\idx)$ and $e=e_\lay(\idx)$:
\[
\begin{aligned}
  (c_\lay,\sigma_\lay)&\gets\OtsSign(P,\lay,\tau,e,M_\lay),\\
  A_\lay&\gets\TreePath(P,\lay,\tau,e),\\
  M_{\lay-1}&\gets\TreeRoot(P,\lay,\tau),
\end{aligned}
\]
returning $\bot$ if $\OtsSign$ does.
\item Output $\sigma=\left(\rho,(s_\kappa,B_\kappa)_{\kappa<k-1},(c_\lay,\sigma_\lay,A_\lay)_{\lay<d}\right)$.
\end{enumerate}
$M_{-1}$ is computed and discarded: it is $\rootnode$ whenever the signer is honest. The signature is serialized as
\[
  \rho
  \concat s_0\concat B_{0,0}\concat\cdots\concat B_{0,a-1}
  \concat\cdots
  \concat s_{k-2}\concat B_{k-2,0}\concat\cdots\concat B_{k-2,a-1}
\]
\[
  \concat\LE_{\lctr}(c_0)\concat\sigma_{0,0}\concat\cdots\concat\sigma_{0,v-1}\concat A_{0,0}\concat\cdots\concat A_{0,h_0-1}
  \concat\cdots
\]
ending with layer $d-1$'s counter, one-time signature and path, of size
\[
  \underbrace{16}_{\rho}+\underbrace{(k-1)(1+a)\cdot 16}_{\text{few-time}}+\underbrace{d\,(4+v\cdot 16)}_{\text{one-time}}+\underbrace{h\cdot 16}_{\text{paths}}
  =16+2464+2028+416=4924\text{ bytes.}
\]

The digest loop resamples $\rho$ and $\OtsSign$ enumerates $c$; both test a random oracle output, admissible with probability $2^{-a}$ and $p=2^{-13.60}$, so each succeeds within its bound except with probability far below $2^{-256}$. On $\bot$ a caller may simply call $\Sig$ again: the digest loop draws fresh randomness, so nothing is permanently unsignable.

\section{Verification}
\label{sec:ver}

$\Ver(\pk,m,\bot)=0$; otherwise $\Ver(\pk,m,\sigma)$ parses $\pk=(\rootnode,P)$ and $\sigma$ as above and:
\begin{enumerate}[leftmargin=2em]
\item Computes $\Digest(P,\rootnode,m,\rho)$, giving $\idx$ and $u$, and rejects unless $u_{k-1}=0$.
\item $M_{d-1}\gets\FtsRec(P,\idx,u,s,B)$.
\item For $\lay=d-1$ down to $0$, with $\tau=\tau_\lay(\idx)$ and $e=e_\lay(\idx)$:
\[
  X\gets\OtsLeaf(P,\lay,\tau,e,M_\lay,c_\lay,\sigma_\lay),\qquad
  M_{\lay-1}\gets\TreeFold(P,\lay,\tau,e,X,A_\lay),
\]
rejecting if $\OtsLeaf$ returns $\bot$.
\item Accepts if and only if $M_{-1}=\rootnode$.
\end{enumerate}

\section{Costs}
\label{sec:costs}

Costs are those of the seed-derived implementation of Remark~\ref{rem:seed}, counted in calls to $\hash$: one per chain step, per Merkle node, per derived secret, and one for each of $\Th$'s other uses whatever its argument length. A one-time key and its leaf therefore cost $v+v(2^w-1)+1=42+294+1=337$.

\begin{center}
\begin{tabular}{@{}lrl@{}}
\toprule
 & Hashes & \\
\midrule
Key generation & 1.38M & layer $0$'s tree, $2^{h_0}\cdot337+2^{h_0}-1$\\
Verification & 497 & \\
\quad $\Digest$ & 1 & \\
\quad $\FtsRec$ & 155 & $(k-1)(1+a)$ nodes $+$ 1 for the roots\\
\quad $\OtsLeaf$ & 315 & $d\left[(2^w-1)v-T+2\right]$, that is $3\times105$\\
\quad $\TreeFold$ & 26 & $h$ nodes\\
Signing & 190K & \\
\quad few-time trees & 43.0K & $(k-1)\left[2\cdot2^a+2^a-1\right]+1$\\
\quad digest attempts & 1.02K & $2^a$ attempts at one hash\\
\quad layer $0$, cached & 21.7K & Remark~\ref{rem:cache}\\
\quad layers 1 and 2 & 86.5K & $2\left[2^{7}\cdot337+2^{7}-1\right]$\\
\quad counter search & 37.3K & $d\lceil1/p\rceil$, $p=2^{-13.60}$\\
\bottomrule
\end{tabular}
\end{center}

\begin{remark}[Test vectors]
Nothing implements this scheme yet, so the document carries no test vectors. One key pair, one message, and the resulting $\rho$, $\idx$, the $k$ indices $u_\kappa$, the $d$ counters $c_\lay$ and the 4924 signature bytes would pin every convention above, and should be added before anyone implements against it.
\end{remark}

\section{Security}
\label{sec:security}

\subsection{Classical security}

\begin{definition}[Strong unforgeability in the ROM]
Let $\SIG=(\Gen,\Sig,\Ver)$ be a signature scheme whose algorithms use a hash function $\hash:\bits{*}\to\bits{256}$. Consider the following game between a signer and an adversary $\mathcal A$ (an arbitrary probabilistic algorithm with unbounded running time and memory), with $\hash$ sampled as a random oracle, that both signer and adversary can query. The signer first runs $(\pk,\sk)\gets\Gen$ and gives $\pk$ to $\mathcal A$. The adversary may then adaptively take any of the following actions:
\begin{enumerate}[leftmargin=2em]
  \item Query the random oracle on any input and receive its 256-bit output.
  \item Submit a message $m\in\bits{\lmsg}$ and receive $\sigma\gets\Sig(\sk,m)$ from the signer, which may be $\bot$. It may do so at most $\qs$ times, on any messages, the same one included: $\Sig$ keeps no state, so nothing here is used up.
  \item Terminate with a claimed forgery $(m^*,\sigma^*)$.
\end{enumerate}
The adversary wins if $\Ver(\pk,m^*,\sigma^*)=1$ and the signer did not return $\sigma^*$ in response to a signing query for $m^*$, meaning:
\begin{itemize}[leftmargin=2em]
  \item if the adversary never queried a signature for $m^*$;
  \item or it did, but no answer it received was $\sigma^*$.
\end{itemize}

Call $\mathcal A$ $q$-bounded if the experiment makes at most $q$ random-oracle queries on every execution, counting those of key generation, signing, and the final verification of the claimed forgery. We say that $\SIG$ has $x$ bits of classical strong unforgeability in the ROM at $\qs$ signatures if every $q\geq1$ and every $q$-bounded $\mathcal A$ satisfy
\[
  \Pr[\mathcal A\text{ wins}]\leq\frac{q}{2^{x}}.
\]
\end{definition}

TODO prove 127 bits of classical strong unforgeability in the ROM at $\qs=2^{24}$ signatures.

\subsection{Quantum security}
\label{sec:quantum}

TODO.

\bibliographystyle{alphaurl}
\bibliography{refs}

\end{document}
